% *======================================================================*
%  Cactus Thorn template for ThornGuide documentation
%  Author: Ian Kelley
%  Date: Sun Jun 02, 2002
%  $Header: /cactusdevcvs/Cactus/doc/ThornGuide/template.tex,v 1.12 2004/01/07 20:12:39 rideout Exp $
%
%  Thorn documentation in the latex file doc/documentation.tex
%  will be included in ThornGuides built with the Cactus make system.
%  The scripts employed by the make system automatically include
%  pages about variables, parameters and scheduling parsed from the
%  relevant thorn CCL files.
%
%  This template contains guidelines which help to assure that your
%  documentation will be correctly added to ThornGuides. More
%  information is available in the Cactus UsersGuide.
%
%  Guidelines:
%   - Do not change anything before the line
%       % START CACTUS THORNGUIDE",
%     except for filling in the title, author, date, etc. fields.
%        - Each of these fields should only be on ONE line.
%        - Author names should be separated with a \\ or a comma.
%   - You can define your own macros, but they must appear after
%     the START CACTUS THORNGUIDE line, and must not redefine standard
%     latex commands.
%   - To avoid name clashes with other thorns, 'labels', 'citations',
%     'references', and 'image' names should conform to the following
%     convention:
%       ARRANGEMENT_THORN_LABEL
%     For example, an image wave.eps in the arrangement CactusWave and
%     thorn WaveToyC should be renamed to CactusWave_WaveToyC_wave.eps
%   - Graphics should only be included using the graphicx package.
%     More specifically, with the "\includegraphics" command.  Do
%     not specify any graphic file extensions in your .tex file. This
%     will allow us to create a PDF version of the ThornGuide
%     via pdflatex.
%   - References should be included with the latex "\bibitem" command.
%   - Use \begin{abstract}...\end{abstract} instead of \abstract{...}
%   - Do not use \appendix, instead include any appendices you need as
%     standard sections.
%   - For the benefit of our Perl scripts, and for future extensions,
%     please use simple latex.
%
% *======================================================================*
%
% Example of including a graphic image:
%    \begin{figure}[ht]
% 	\begin{center}
%    	   \includegraphics[width=6cm]{MyArrangement_MyThorn_MyFigure}
% 	\end{center}
% 	\caption{Illustration of this and that}
% 	\label{MyArrangement_MyThorn_MyLabel}
%    \end{figure}
%
% Example of using a label:
%   \label{MyArrangement_MyThorn_MyLabel}
%
% Example of a citation:
%    \cite{MyArrangement_MyThorn_Author99}
%
% Example of including a reference
%   \bibitem{MyArrangement_MyThorn_Author99}
%   {J. Author, {\em The Title of the Book, Journal, or periodical}, 1 (1999),
%   1--16. {\tt http://www.nowhere.com/}}
%
% *======================================================================*

% If you are using CVS use this line to give version information
% $Header: /cactusdevcvs/Cactus/doc/ThornGuide/template.tex,v 1.12 2004/01/07 20:12:39 rideout Exp $

\documentclass{article}

% Use the Cactus ThornGuide style file
% (Automatically used from Cactus distribution, if you have a
%  thorn without the Cactus Flesh download this from the Cactus
%  homepage at www.cactuscode.org)
\usepackage{../../../../doc/latex/cactus}

\begin{document}

% The author of the documentation
\author{Roland Haas \textless roland.haas@physics.gatech.edu\textgreater}

% The title of the document (not necessarily the name of the Thorn)
\title{SphereRad}

% the date your document was last changed, if your document is in CVS,
% please use:
%    \date{$ $Date: 2004/01/07 20:12:39 $ $}
\date{Tue Mar 24 12:14:08 EDT 2009}

\maketitle

% Do not delete next line
% START CACTUS THORNGUIDE

% Add all definitions used in this documentation here
%   \def\mydef etc
\newcommand{\thorn}[1]{\texttt{#1}}

% Add an abstract for this thorn's documentation
\begin{abstract}
    SphereRad is a utility thorn that can set the radius and position of
    spherical surfaces (see thorn \thorn{SphericalSurface}) from a number of
    sources.
\end{abstract}

% The following sections are suggestive only.
% Remove them or add your own.

\section{Introduction}
A number of thorns operate on spherical surfaces as 
repositories to store location and radius of a topological sphere (e.g.\
\thorn{ShiftTracker}, \thorn{MinTracker}) while others use these surfaces as
input (e.g.\
\thorn{IHSpin}). This thorn serves as an interface to either add radius information
for setter thorns that do not provide it (\thorn{MinTracker}, \thorn{ShiftTracker}) and can
also copy data from \thorn{ShiftTracker} to the spherical surface. Finally it allows
for specification of location and radius via parameters so that a fixed
surface can be set up. It only support the creation of coordinate spheres i.e.\ .
spheres of constant coordinate radius.

\section{Using This Thorn}\label{sec:sphererad_usingthisthorn}
Use the parameter \verb|surface_index| to select the target spherical surface.
If the source is \thorn{ShiftTracker} then set the corresponding entry in
\verb|shifttracker_index| to indicate which tracker to use, the radius of the
surface will be taken from \verb|radius_surface|. If the source is another
spherical surface (e.g.\ . set by \thorn{MinTracker}) then set \verb|sperical_surface| to
\verb|true| in which case \verb|shifttracker_index| is interpreted as a
spherical surface number, not as a tracker number. The radius is still taken
from \verb|radius_surface|, not from the source spherical surface's
\verb|sf_radius| values. Finally if \verb|force_position| is set, then the
centre coordinates are taken from the parameters \verb|x_fixed\|,
\verb|y_fixed\|, and \verb|z_fixed\| respectively.

SphereRad can do a simply centre of mass calculation for surfaces for which
\verb|center_of_mass| is \verb|true|. In this case the centre coordinate of
the spherical surface is calculated not from any parameter but is set to the
centre of mass of surfaces $0$ and $1$ only. To this end the parameters
\verb|mass[0]| and \verb|mass[1]| are used for the masses of these two
spherical surfaces, other entries in \verb|mass| are ignored. With these the
centre of mass is calculated as
\begin{verbatim}
  const CCTK_REAL total_mass=mass[0]+mass[1];
  const CCTK_REAL xcm=(xah[0]*mass[0]+xah[1]*mass[1]) /total_mass;
  const CCTK_REAL ycm=(yah[0]*mass[0]+yah[1]*mass[1]) /total_mass;
  const CCTK_REAL zcm=(zah[0]*mass[0]+zah[1]*mass[1]) /total_mass; .
\end{verbatim}

\subsection{Obtaining This Thorn}
This thorns lives in the PSUThorns repository.

\subsection{Basic Usage}
For basic use, use \thorn{ShiftTracker} to track the puncture(s) and set
\verb|shifttracker_index| to the correct values. Set \verb|radius_surface| to
a value outside of the apparent horizon, but still far away from any other
horizons. A good value might be of the order of the puncture mass $M_{BH}$.
See sec.~\ref{sec:sphererad_usingthisthorn} for a more thorough description of
the thorn.

\subsection{Special Behaviour}
\begin{itemize}
\item ignores \verb|mass[n]| with $n \ge 2$.
\item treats \verb|shifttracker_index| as a spherical surface index if
    \verb|spherical_surface| is set.
\end{itemize}

\subsection{Interaction With Other Thorns}
Can read out \thorn{ShiftTracker}'s variables to set the spherical surface. Cannot
read out \thorn{MinTracker}s variables.

\subsection{Examples}
Taken from R1.rpar:
\begin{verbatim}
#############################################################
# SphereRad
#############################################################

SphereRad::max_nsurface           = 3
SphereRad::surface_index[0]       = 2
SphereRad::surface_index[1]       = 3
SphereRad::surface_index[2]       = 4
SphereRad::shifttracker_index[0]  = 0
SphereRad::shifttracker_index[1]  = 1
SphereRad::force_position[2]      = "yes"
# surface no. 2 defaults to centre at the origin
SphereRad::radius_surface[0]      = 0.4
SphereRad::radius_surface[1]      = 0.4
# BH bare masses: 0.483
SphereRad::radius_surface[2]      = 1
SphereRad::verbose                = 0

#############################################################
# Shift tracker
#############################################################

ShiftTracker::x0[0]              = $xp
ShiftTracker::y0[0]              = 0.0
ShiftTracker::z0[0]              = 0.0
#ShiftTracker::surface_index[0]   = 0

ShiftTracker::x0[1]              = $xm
ShiftTracker::y0[1]              = 0.0
ShiftTracker::z0[1]              = 0.0
#ShiftTracker::surface_index[1]   = 1

ShiftTracker::num_trackers       = 2
ShiftTracker::verbose            = 0
ShiftTracker::output_every       = 1
ShiftTracker::interpolator_name  = "Lagrange polynomial interpolation"
ShiftTracker::interpolator_pars  = "order=4"
ShiftTracker::beta1_var          = "Kranc2BSSNChi::beta1"
ShiftTracker::beta2_var          = "Kranc2BSSNChi::beta2"
ShiftTracker::beta3_var          = "Kranc2BSSNChi::beta3"
\end{verbatim}

\subsection{Support and Feedback}
very little I'm afraid.

\section{History (darcs changes)}
\begin{verbatim}
Wed Nov  7 15:09:37 EST 2007  herrmann@gravity.psu.edu
  * SphereRad: handle sf_active and test for maximum number of surfaces

Thu Jun 21 15:37:07 EDT 2007  herrmann@gravity.psu.edu
  * SphereRad: more max. no. of spheres (10->100)

Fri Feb 16 16:43:59 EST 2007  herrmann@gravity.psu.edu
  * SphereRad: center of mass computation

  can now define a surface tracking the center of mass.

Fri Feb  2 01:02:30 EST 2007  herrmann@gravity.psu.edu
  * SphereRad: more verbose output

Wed Jan 31 16:53:34 EST 2007  herrmann@gravity.psu.edu
  * SphereRad: thorn to set spherical surface

  sets spherical surface of constant radius. used to feed IHSpin.
\end{verbatim}

\subsection{Acknowledgements}


\begin{thebibliography}{9}

\end{thebibliography}

% Do not delete next line
% END CACTUS THORNGUIDE

\end{document}
